% MÓVIL - fichas individuales de ejecución (27 casos: CP-055..CP-081)
% GUÍA EQUIPO (Carlos Ramos): completar los parámetros de \casoejec{ID}{título}{resultado obtenido}{PASS|FAIL|BLOCKED|N/A}{evidencia o log}{DEF-XX}.
% - Estado: solo PASS, FAIL, BLOCKED o N/A.
% - Defecto: ID del defecto (ej. DEF-01) si falló, o Ninguno si pasó.
% - Resultado obtenido: describir lo observado al ejecutar la prueba.
% - Si existe evidencias/<ID>.png, la imagen se inserta automáticamente dentro de su ficha.
\par\medskip\noindent\textbf{Módulo Móvil --- Seguridad y operación de entregas (27 casos: CP-055 a CP-081).}\par\smallskip
En esta sección se presentan los resultados de la ejecución de los 27 casos de prueba del módulo
móvil (\texttt{CP-055} a \texttt{CP-081}): control de acceso por rol y por recurso, planificación del
día operativo, urgencias, entregas fallidas y su gestión posterior, recorrido de extremo a extremo y
cierre de ruta con regreso a tienda.

Todas las pruebas se ejecutaron en local. El backend NestJS (commit \texttt{9b17cce}) se levantó en
el puerto \texttt{3003} sobre una base PostgreSQL~16 aislada creada solo para esta ejecución
(\texttt{logistica\_ape4\_movil}, 22 migraciones aplicadas y datos demo sintéticos), con el simulador
de pagos habilitado. La aplicación Flutter (\texttt{app-movil-logistica}, commit \texttt{89f7050}) se
compiló en modo \emph{debug} y se instaló en un teléfono físico Xiaomi M2101K6G con Android~13
conectado por cable USB; el teléfono alcanzó la API local mediante \texttt{adb reverse tcp:3003}.

Los casos de API se automatizaron con el script \texttt{evidencias/scripts/ape4-movil.mjs}, que inicia
sesión con cuentas reales (\texttt{ADMIN}, \texttt{CLIENT-A}, \texttt{CLIENT-B}, \texttt{DRIVER-A},
\texttt{DRIVER-B}) y llama a la API por HTTP con sus tokens JWT. Cada caso guarda en
\texttt{evidencias/json/CP-XXX.json} las peticiones, las respuestas y cada verificación del criterio de
aceptación; cuando el criterio exige comprobar persistencia (reservas, cuentas, retiros) se consultó
además la base de datos. En los casos de rechazo se comparó el estado antes y después de cada intento:
un caso solo se consideró exitoso si, además del código HTTP esperado, no dejó efectos parciales. Las
ausencias registradas durante las pruebas se anularon al terminar cada caso.

Los casos que involucran la pantalla del repartidor (\texttt{CP-057}, \texttt{CP-060},
\texttt{CP-067}, \texttt{CP-073} a \texttt{CP-077} y \texttt{CP-079}) se repitieron en el teléfono con
el script de apoyo \texttt{evidencias/scripts/ape4-movil-telefono.mjs}, que registra el estado real
del backend después de cada acción en la app (\texttt{evidencias/json/movil-telefono.json}). La
ejecución en el teléfono reveló dos defectos de la app: no carga rutas creadas sin zona
(\texttt{DEF-06}) y no ofrece las acciones de entrega fallida, cierre de ruta ni confirmación de
regreso (\texttt{DEF-07}). Ambos se detallan en las fichas y en el registro de incidentes.

% ---------------------------------------------------------------------------
\par\medskip\noindent\textbf{Seguridad y control de acceso (CP-055 a CP-060).}\par\smallskip
Se comprobó que cada actor solo accede a sus propios recursos y que ningún rechazo filtra datos ni
modifica el estado.

\casoejec{CP-055}{Acceso sin sesión JWT a rutas, pedidos y configuración}{Las tres peticiones sin
cabecera \texttt{Authorization} (\texttt{GET /drivers/me/routes}, \texttt{GET /shop/orders} y
\texttt{GET /admin/shop-settings}) respondieron HTTP 401. Ninguna respuesta incluyó listados, rutas,
paradas ni parámetros de configuración.}{PASS}{Script \texttt{ape4-movil.mjs};
JSON: \texttt{evidencias/json/CP-055.json}}{Ninguno}

\casoejec{CP-056}{Aislamiento entre CLIENT-B y pedido/\allowbreak{}comprobante de CLIENT-A}{Con el JWT de
CLIENT-B sobre una compra Shop de CLIENT-A: \texttt{GET /shop/orders/:id} y \texttt{/receipt}
respondieron 404, y \texttt{GET /orders/:id} y \texttt{PATCH /orders/:id/status} respondieron 403.
Ninguna respuesta contenía el destinatario, la dirección, el número comercial ni el total de CLIENT-A,
y el pedido consultado como ADMIN quedó idéntico antes y después.}{PASS}{JSON:
\texttt{evidencias/json/CP-056.json}}{Ninguno}

\casoejec{CP-057}{Aislamiento de DRIVER-B ante ruta y parada de DRIVER-A}{\emph{API:} Mi Ruta de
DRIVER-B no incluyó la ruta de DRIVER-A; consultarla (\texttt{GET /routes/:id}), iniciarla, registrar
un fallo en su parada y cerrarla como DRIVER-B respondieron 403, y la ruta y sus paradas conservaron el
mismo estado. \emph{Teléfono:} con la sesión del segundo repartidor la app mostró ``No tienes rutas
para hoy'', sin listar la ruta asignada al primer repartidor.}{PASS}{JSON:
\texttt{evidencias/json/CP-057.json}; captura del teléfono}{Ninguno}

\casoejec{CP-058}{Registro público con rol ADMIN rechazado sin escalación}{\texttt{POST
/auth/register} con un correo nuevo y el campo adicional \texttt{role: ADMIN} respondió 400. El inicio
de sesión posterior con esas credenciales respondió 401. En la base de datos no existe ninguna cuenta
con ese correo y el número de cuentas ADMIN siguió siendo 1 (la del seed).}{PASS}{JSON y consulta a
la base: \texttt{evidencias/json/CP-058.json}}{Ninguno}

\casoejec{CP-059}{Intento de CLIENT de modificar configuración o cobertura}{Como CLIENT-A,
\texttt{PATCH /admin/shop-settings} y \texttt{PUT /admin/store-coverage} (nuevo polígono) respondieron
403. La configuración de la tienda y la cobertura (ubicación y polígono) leídas como ADMIN fueron
idénticas antes y después.}{PASS}{JSON comparativo: \texttt{evidencias/json/CP-059.json}}{Ninguno}

\casoejec{CP-060}{Acceso autorizado por titular, ADMIN y DRIVER sobre recursos propios}{CLIENT-A
leyó su pedido y su comprobante (200); su historial devolvió 3 pedidos, todos propios. ADMIN consultó
la configuración y la bandeja de despacho (200). DRIVER-A obtuvo 200 en Mi Ruta solo con su ruta
asignada (0 rutas de otros repartidores). En el teléfono, el repartidor vio su ruta del día (zona
Centro, Asignada, 0 de 2 entregados); esta prueba se hizo con una ruta creada con zona, por el defecto
\texttt{DEF-06} descrito en \texttt{CP-073}.}{PASS}{JSON: \texttt{evidencias/json/CP-060.json}; captura
del teléfono}{Ninguno}

% ---------------------------------------------------------------------------
\par\medskip\noindent\textbf{Planificación del día operativo y urgencias (CP-061 a CP-066).}\par\smallskip
Se verificó que solo los pedidos del día operativo (America/Guayaquil) se conviertan en rutas y que
las urgencias nunca alteren una ruta iniciada sin decisión del ADMIN.

\casoejec{CP-061}{Pedido de hoy elegible en propuesta, ruta y asignación}{Con dos pedidos para hoy
(16:00 y 16:10), \texttt{GET /route-proposals} respondió 200 con \texttt{isToday=true} e incluyó
ambos. \texttt{POST /routes} creó la ruta en \texttt{DRAFT} (201) y \texttt{POST
/routes/:id/assign} la dejó \texttt{ASSIGNED} a DRIVER-A (200), con 2 paradas y la fecha de
hoy.}{PASS}{JSON: \texttt{evidencias/json/CP-061.json}}{Ninguno}

\casoejec{CP-062}{Pedido futuro admitido en compra y bloqueado en ruta/\allowbreak{}assign}{La compra
Shop y el pedido ADMIN para mañana se crearon (201). La vista previa de mañana respondió 200 con
\texttt{isToday=false} y \texttt{canConfirm=false} e incluyó ambos pedidos. \texttt{POST /routes} para
mañana respondió 409 \texttt{ROUTE\_DATE\_NOT\_TODAY} y la asignación directa \texttt{POST
/routes/assign} también 409. No se creó ninguna ruta y los dos pedidos siguieron
\texttt{PENDING}.}{PASS}{JSON: \texttt{evidencias/json/CP-062.json}}{Ninguno}

\casoejec{CP-063}{Ausencia registrada hoy impide asignar repartidor}{Se registró una ausencia de hoy
(201) a un repartidor sin rutas, para que la ausencia fuera la única causa posible. Su disponibilidad
devolvió \texttt{isAvailable=false} con la ausencia indicada y el intento de asignarle una ruta
respondió 409. La ruta quedó en \texttt{DRAFT} sin conductor y no se creó ninguna reserva (0
reservas).}{PASS}{JSON: \texttt{evidencias/json/CP-063.json}}{Ninguno}

\casoejec{CP-064}{Ausencia futura no altera ni reasigna rutas activas de hoy}{Con la ruta de hoy de
DRIVER-A asignada, se registró una ausencia de mañana a pasado mañana (201, \texttt{ACTIVE}). La ruta
mantuvo el mismo conductor, estado, paradas e intentos y su reserva siguió activa (1 reserva): no hubo
reasignación, cancelación ni liberación automática. La ausencia se anuló al terminar.}{PASS}{JSON
comparativo: \texttt{evidencias/json/CP-064.json}}{Ninguno}

\casoejec{CP-065}{Pedido urgente con ruta IN\_PROGRESS no altera ruta iniciada}{Con la ruta de
DRIVER-A ya iniciada (\texttt{IN\_PROGRESS}), un pedido urgente creado para hoy obtuvo la propuesta
\texttt{NEW\_ROUTE} en \texttt{GET /dispatch/orders/:id/proposal}, sin proponer la ruta iniciada. Las
paradas, el orden y el estado de la ruta iniciada fueron idénticos antes y después.}{PASS}{JSON
comparativo: \texttt{evidencias/json/CP-065.json}}{Ninguno}

\casoejec{CP-066}{Decisión manual ADMIN en urgencias: INSERT con cupo o NEW\_ROUTE}{Mientras había
rutas no iniciadas con cupo, cada urgente obtuvo la propuesta \texttt{INSERT}. La ruta destino no
cambió hasta que el ADMIN confirmó explícitamente (\texttt{CONFIRM}, 200 \texttt{APPLIED}); cada
confirmación agregó exactamente una parada con secuencia contigua. Se hicieron 5 inserciones en dos
rutas hasta llenarlas (4 paradas cada una). El siguiente urgente obtuvo \texttt{NEW\_ROUTE} y, al
confirmarlo con un repartidor libre, se creó una ruta nueva. Ninguna ruta superó 4 paradas, las
decisiones quedaron auditadas con el actor ADMIN y la ruta iniciada no recibió
inserciones.}{PASS}{JSON y auditoría de despacho: \texttt{evidencias/json/CP-066.json}}{Ninguno}

% ---------------------------------------------------------------------------
\par\medskip\noindent\textbf{Entrega fallida y gestión posterior (CP-067 a CP-071).}\par\smallskip
Se verificó que una visita fallida se conserve en el historial, que solo el repartidor titular pueda
registrarla y que las gestiones del ADMIN (reprogramar o retiro en tienda) no dupliquen entregas.

\casoejec{CP-067}{Entrega fallida en móvil conserva parada e historial pendiente}{\emph{API:} DRIVER-A
registró la entrega fallida con motivo \texttt{RECIPIENT\_ABSENT} (200). La parada pasó a
\texttt{FAILED} conservando el intento con su motivo, el pedido no quedó \texttt{DELIVERED} ni
\texttt{CANCELLED} y el seguimiento del cliente lo mostró pendiente de gestión.
\emph{Teléfono:} la app no permite registrar el fallo; la parada en curso solo ofrece ``En camino'' y
``Entregado''. El fallo tuvo que registrarse por API y la app solo lo mostró después (``Intento
fallido''). El paso ``registrar entrega fallida desde móvil'' no se puede ejecutar.}{FAIL}{JSON:
\texttt{evidencias/json/CP-067.json} y \texttt{movil-telefono.json}; captura del teléfono}{DEF-07}

\casoejec{CP-068}{Registro de fallo por actor o conductor no autorizado rechazado}{El intento de
registrar un fallo sobre la parada de DRIVER-A respondió 403 tanto con el JWT de DRIVER-B como con el
de CLIENT-A. La parada siguió \texttt{PENDING} sin intentos y el pedido no cambió de
estado.}{PASS}{JSON: \texttt{evidencias/json/CP-068.json}}{Ninguno}

\casoejec{CP-069}{Gestión post-fallo ADMIN: CONTACT y RESCHEDULE con nueva fecha}{Tras una visita
fallida y la ruta cerrada, \texttt{CONTACT} con acuerdo explícito (mañana 12:00, ventana 11:30--12:30)
y \texttt{RESCHEDULE} con ese acuerdo respondieron 200 \texttt{APPLIED}. El pedido quedó
\texttt{PENDING}, sin ruta y con la nueva fecha (\texttt{2026-11-15T17:00:00Z}). La ruta original
conservó la parada \texttt{FAILED} con su intento, no quedó ninguna asignación activa duplicada y el
pedido apareció en la propuesta de mañana.}{PASS}{JSON: \texttt{evidencias/json/CP-069.json}}{Ninguno}

\casoejec{CP-070}{Gestión post-fallo ADMIN: acuerdo de retiro en tienda (PICKUP\_AGREE)}{\texttt{PICKUP\_AGREE}
respondió 200 \texttt{APPLIED}. El pedido quedó \texttt{PENDING} con el retiro en
\texttt{PENDING\_PICKUP} (sin \texttt{pickedUpAt}), sin marcarse entregado. El seguimiento del cliente
mostró el hito ``Pendiente de retiro en tienda'', la visita fallida se conservó y el stock no cambió
(104 $\rightarrow$ 104).}{PASS}{JSON: \texttt{evidencias/json/CP-070.json}}{Ninguno}

\casoejec{CP-071}{Confirmación de retiro (PICKUP\_CONFIRM) y rechazo de repetición}{La primera
\texttt{PICKUP\_CONFIRM} respondió 200: el pedido quedó \texttt{DELIVERED} con el retiro en
\texttt{PICKED\_UP}. La repetición se rechazó con 409, pero con el mensaje genérico ``Commercially
resolved order cannot be managed for pickup'' y sin el código \texttt{PICKUP\_ALREADY\_CONFIRMED} que
exige el caso. La integridad sí se cumplió: un solo registro de retiro y stock sin cambio (104
$\rightarrow$ 104).}{FAIL}{JSON: \texttt{evidencias/json/CP-071.json}}{DEF-05}

% ---------------------------------------------------------------------------
\par\medskip\noindent\textbf{Recorrido de extremo a extremo (CP-072 a CP-075).}\par\smallskip
Se recorrió el flujo completo, desde la compra del cliente hasta la entrega en el teléfono del
repartidor, contrastando lo que muestra la app con lo que registra el servidor.

\casoejec{CP-072}{Recorrido E2E: creación de compra y verificación de estado inicial}{CLIENT-A compró 2
unidades con la tarjeta demo aprobada (201). El pedido se creó en \texttt{PENDING} con importes
congelados: subtotal 5180, IVA 15\,\% 777 y total 5957 centavos, igual a lo cotizado. El seguimiento y
el comprobante de CLIENT-A mostraron el mismo pedido y total, y el stock bajó de 109 a 107.}{PASS}{JSON:
\texttt{evidencias/json/CP-072.json}}{Ninguno}

\casoejec{CP-073}{Recorrido E2E: propuesta, confirmación y asignación en día operativo}{\emph{API:} la
propuesta de hoy incluyó el pedido; el borrador se persistió (201, \texttt{DRAFT}) y quedó
\texttt{ASSIGNED} a DRIVER-A, con una sola reserva y una sola aparición en \texttt{GET
/drivers/me/routes}. \emph{Teléfono:} la ruta confirmada desde la propuesta no tiene zona
(\texttt{"zone": null}) y la app no la muestra. Mi Ruta respondió ``No se pudieron cargar tus rutas. El
servidor devolvió rutas inválidas''. La app exige que \texttt{zone} sea un objeto, y con una sola ruta
sin zona el repartidor no ve ninguna de sus rutas. Con una ruta creada con zona, la app la listó
correctamente.}{FAIL}{JSON: \texttt{evidencias/json/CP-073.json}; captura del
teléfono}{DEF-06}

\casoejec{CP-074}{Recorrido E2E: progreso móvil reflejado en seguimiento del cliente}{\emph{API:} con
la compra de CP-072, cada transición (inicio, \texttt{IN\_TRANSIT}, \texttt{DELIVERED}) se reflejó
igual en el backend y en el seguimiento de CLIENT-A, sin coordenadas GPS en la línea de tiempo.
\emph{Teléfono:} el repartidor inició la ruta y marcó ``En camino'' y ``Entregado'' en la parada de una
compra Shop. Tras cada acción, el seguimiento del cliente coincidió con el backend: \texttt{IN\_TRANSIT}
y luego \texttt{DELIVERED}, con la línea de tiempo PENDING $\rightarrow$ ASSIGNED $\rightarrow$
IN\_TRANSIT $\rightarrow$ DELIVERED.}{PASS}{JSON: \texttt{evidencias/json/CP-074.json} y
\texttt{movil-telefono.json}; captura del teléfono}{Ninguno}

\casoejec{CP-075}{Recorrido E2E: desconexión móvil y recuperación al reconectar}{En el teléfono se
cortó la conexión con el servidor (se quitó \texttt{adb reverse}) y se pulsó ``Iniciar ruta''. La app
no fingió éxito: mostró ``No se pudo conectar con el servidor. Actualiza la lista antes de volver a
intentar iniciar una ruta'', la ruta siguió ``Asignada'' y en el backend continuó \texttt{ASSIGNED} con
\texttt{startedAt} nulo. Al reconectar y actualizar, la app mostró el estado real del servidor, sin
cola offline. El inicio posterior se registró una sola vez.}{PASS}{JSON:
\texttt{evidencias/json/CP-075.json}; captura del teléfono}{Ninguno}

% ---------------------------------------------------------------------------
\par\medskip\noindent\textbf{Cierre de ruta y regreso a tienda (CP-076 a CP-081).}\par\smallskip
Se verificó que la ruta solo se cierre con todas las visitas resueltas y que el repartidor quede
ocupado hasta confirmar su regreso físico a la tienda.

\casoejec{CP-076}{Cierre de ruta móvil rechazado con visita aún PENDING}{Con una parada todavía
\texttt{PENDING}, \texttt{POST /routes/:id/complete} de DRIVER-A respondió 409. La ruta siguió
\texttt{IN\_PROGRESS} sin \texttt{completedAt} y el repartidor siguió ocupado. En el teléfono se
repitió el escenario (parada 1 con intento fallido y parada 2 pendiente): el backend rechazó el cierre
con 409 (``One to four resolved verified visits required'') y la app mostró la ruta ``En curso'' con la
parada pendiente. La app no tiene botón de cierre (\texttt{DEF-07}).}{PASS}{JSON:
\texttt{evidencias/json/CP-076.json}; captura del teléfono}{Ninguno}

\casoejec{CP-077}{Cierre de ruta COMPLETED con visitas resueltas y una FAILED}{DRIVER-A entregó la
segunda parada y cerró la ruta: 200 \texttt{COMPLETED} con \texttt{completedAt}. La parada fallida
siguió \texttt{FAILED}, su pedido no se marcó \texttt{DELIVERED} y el repartidor siguió no disponible
con la ruta como bloqueante hasta su regreso. En el teléfono, después del cierre (hecho por API con el
token del repartidor, ver \texttt{DEF-07}), la app mostró ``Completada · 1 de 2 entregados · 1 con
intento fallido · Sigues ocupado: regresa al depósito y confirma tu llegada''.}{PASS}{JSON:
\texttt{evidencias/json/CP-077.json} y \texttt{movil-telefono.json}; captura del teléfono}{Ninguno}

\casoejec{CP-078}{Bloqueo de nueva salida antes de confirmar regreso (returnedAt)}{Con la ruta previa
\texttt{COMPLETED} y \texttt{returnedAt} nulo, asignar a DRIVER-A una nueva ruta de hoy respondió 409.
La segunda ruta quedó en \texttt{DRAFT} sin conductor y DRIVER-A conservó una sola reserva activa, sin
doble reserva.}{PASS}{JSON: \texttt{evidencias/json/CP-078.json}}{Ninguno}

\casoejec{CP-079}{Confirmación de regreso por titular libera y actualiza disponibilidad}{\texttt{POST
/routes/:id/return} de DRIVER-A respondió 200 con \texttt{returnedAt}, liberó la reserva (0 activas) y
dejó al repartidor disponible. Con otro repartidor con ausencia vigente hoy, el regreso también se
registró y liberó la reserva, pero su disponibilidad siguió en \texttt{false} por la ausencia. En el
teléfono, tras confirmar el regreso por API (\texttt{DEF-07}), la app mostró ``Visitas finalizadas.
Regreso al depósito confirmado'' y el repartidor quedó disponible.}{PASS}{JSON:
\texttt{evidencias/json/CP-079.json} y \texttt{movil-telefono.json}; captura del teléfono}{Ninguno}

\casoejec{CP-080}{Repetición de regreso a tienda sin duplicar ni liberar dos veces}{Repetir
\texttt{POST /routes/:id/return} respondió 200 con exactamente la misma respuesta que la primera vez
(idempotente). \texttt{returnedAt} y \texttt{returnedById} no cambiaron, no se creó ni liberó otra
reserva y la disponibilidad del repartidor siguió igual.}{PASS}{JSON y consulta a la base:
\texttt{evidencias/json/CP-080.json}}{Ninguno}

\casoejec{CP-081}{Intento de cierre o regreso por DRIVER-B rechazado sobre ruta ajena}{Con el JWT de
DRIVER-B, \texttt{POST /routes/:id/complete} y \texttt{POST /routes/:id/return} sobre la ruta de
DRIVER-A respondieron 403. El estado, \texttt{completedAt}, \texttt{returnedAt} y la reserva activa
de DRIVER-A quedaron exactamente iguales.}{PASS}{JSON: \texttt{evidencias/json/CP-081.json}}{Ninguno}
